\documentclass[11pt]{article}

\usepackage{amsmath,amssymb,amsthm}
\usepackage{bm}
\usepackage{mathtools}
\usepackage[margin=1in]{geometry}
\usepackage{setspace}
\onehalfspacing
\usepackage{enumitem}
\usepackage{booktabs}
\usepackage{hyperref}
\usepackage{algorithm}
\usepackage{algpseudocode}
\usepackage{float}

\newcommand{\R}{\mathbb{R}}
\newcommand{\ind}{\mathbf{1}}
\newcommand{\norm}[1]{\left\lVert #1 \right\rVert}
\newcommand{\trans}{^{\top}}
\newcommand{\logit}{\mathrm{logit}}
\newcommand{\Cat}{\mathrm{Cat}}
\newcommand{\Dir}{\mathrm{Dir}}
\newcommand{\IG}{\mathrm{IG}}
\newcommand{\IW}{\mathcal{IW}}
\newcommand{\tr}{\mathrm{tr}}

\title{\textbf{model\_v14}\\
Bayesian Joint Multilayered LSIRM with a Mixture-of-Mixtures\\
Latent-Position Clustering Prior}
\author{}
\date{May 2026}

\begin{document}
\maketitle

\section{Overview}
\label{sec:overview}

This document specifies \emph{model\_v14}, a Bayesian joint multilayered Latent
Space Item Response Model (LSIRM) for $L=4$ heterogeneous response layers
(binary, continuous, count, and ordinal), augmented with a
\emph{Mixture-of-Mixtures} (MoM) prior on the item latent positions. It
extends \emph{model\_v13}, which used a single-level (over-fitted)
finite Gaussian mixture with sparse Dirichlet weights to cluster the global
item positions $z_q$. The LSIRM likelihood layer, the latent geometry, the
respondent-side parameters, and the rotation/reflection invariance discussion
are identical to v13. The novelty in v14 is the prior on $\{z_q\}$.

\paragraph{Why a Mixture of Mixtures?}
A single Gaussian mixture conflates two notions that are conceptually
distinct in latent-space models:
\begin{enumerate}[label=(\roman*)]
\item a \emph{density basis} --- the Gaussian components needed to approximate
       a possibly non-Gaussian density of latent item positions, and
\item an \emph{interpretable clustering unit} --- contiguous regions of
       latent space that correspond to substantively meaningful groups of
       items (e.g.\ topical or methodological clusters).
\end{enumerate}
A flat Gaussian mixture forces a single label to do both jobs. With many
components, every cluster becomes a single Gaussian, often producing
over-segmentation: a single elongated or skewed cluster gets split into
several Gaussian components, none of which is itself the cluster the user
wants. With few components, the density fit is poor and items are forced
into ill-fitting Gaussians.

The Mixture-of-Mixtures resolves this by introducing a \emph{two-level}
hierarchy:
\begin{itemize}
\item \emph{Component level} ($c_q$, label $l$): provides a flexible
      density basis, one Gaussian per component.
\item \emph{Cluster level} ($g_q$, label $g$): groups components into
      interpretable clusters. Each cluster is itself a finite Gaussian
      mixture, and so can represent non-Gaussian, multi-modal, or
      heavy-tailed cluster shapes.
\end{itemize}
Components are \emph{not} shared across clusters: each cluster $g$ has its
own component pool of size $K_g$. The prior is constructed so that the
components belonging to a single cluster lie tightly together in latent
space (governed by a within-cluster scale $\Sigma_l/\kappa_0$), while
distinct cluster centres $\mu_g^*$ are spread widely (governed by a
between-cluster scale $M_0$). This \emph{scale separation}
$M_0 \gg \tr(\mathbb{E}[\Sigma_l])/\kappa_0$ is what gives the cluster
labels $g_q$ their identifiability: a non-Gaussian cluster can be assembled
from many close-together Gaussians, but components that are far apart
cannot be glued into a single cluster.

The remainder of the document follows the v13 structure: indexing and
notation (Sec.~\ref{sec:indexing}), the LSIRM likelihood
(Sec.~\ref{sec:lsirm}), the new MoM prior (Sec.~\ref{sec:mom-prior},
the principal new section), the remaining priors and rotation invariance
(Secs.~\ref{sec:other-priors}--\ref{sec:rotation}), the joint posterior
(Sec.~\ref{sec:joint}), conditional updates including the new
two-level allocation moves (Secs.~\ref{sec:niw}--\ref{sec:wishart}),
the cluster-level Jain--Neal split--merge derivation
(Sec.~\ref{sec:splitmerge}), posterior summaries
(Sec.~\ref{sec:summaries}), and the full sweep
(Sec.~\ref{sec:sweep}).

\paragraph{Implementation note.}
We continue to assume $L=4$ heterogeneous layers in the order
$\ell=1$ (binary), $\ell=2$ (continuous, Student-$t$ via normal--gamma
scale mixture), $\ell=3$ (count, Negative-Binomial with item-specific
overdispersion), $\ell=4$ (ordinal, LSGRM cumulative logit). The MoM
prior is on the global pool $\{z_q\}_{q=1}^P$ of all item positions
across all layers, exactly as in v13.

\paragraph{Common-respondent assumption.}
As in v13, all $L$ layers share a common respondent set indexed by
$i=1,\ldots,N$ with shared latent positions $a_i\in\R^d$ ($d=2$).
Each layer has its own item set indexed by $j=1,\ldots,J_\ell$ with
layer-specific positions $b_j^{(\ell)}\in\R^d$, and the global item
position vector $z_q$, $q=1,\ldots,P=\sum_\ell J_\ell$, indexes all
$(\ell,j)$ pairs jointly.

\section{Indexing, dimensions, and observed data}
\label{sec:indexing}

Table~\ref{tab:notation} summarises the indices, dimensions, and key
quantities. New v14 quantities relative to v13 are marked with
$^{\dagger}$.

\begin{table}[H]
\centering
\small
\begin{tabular}{l l l}
\toprule
Symbol & Range / type & Meaning \\
\midrule
$i$        & $1,\ldots,N$            & respondent index \\
$\ell$     & $1,\ldots,L=4$          & response-layer index \\
$j$        & $1,\ldots,J_\ell$       & local item index within layer $\ell$ \\
$q$        & $1,\ldots,P=\sum_\ell J_\ell$ & global item index across all layers \\
$d$        & $2$                     & latent-space dimension \\
$a_i$      & $\R^d$                  & respondent latent position \\
$b_j^{(\ell)}$ & $\R^d$              & item latent position (layer $\ell$, local $j$) \\
$z_q$      & $\R^d$                  & global item position, $z_q = b_{j(q)}^{(\ell(q))}$ \\
$\alpha_i^{(\ell)}$ & $\R$           & respondent ability (layer $\ell$) \\
$\beta_j^{(\ell)}$  & $\R$           & item easiness/intercept (layer $\ell$) \\
$\gamma_\ell$ & $\R_{>0}$            & distance scaling (layer $\ell$) \\
$\sigma_{\alpha,\ell}^2$ & $\R_{>0}$ & ability variance (layer $\ell$) \\
$\sigma_0^2$ & $\R_{>0}$             & continuous-layer scale (Student-$t$) \\
$\nu_t$ & $\R_{>0}$                  & Student-$t$ degrees of freedom \\
$\lambda_{ij}^{(\ell)}$ & $\R_{>0}$  & latent precision (continuous layer scale mixture) \\
$\kappa_j$ & $\R_{>0}$               & Negative-Binomial overdispersion (count layer) \\
$\tau_h^{(\ell)}$ & $\R$              & ordinal threshold $h=1,\ldots,H_\ell-1$ \\
\midrule
$c_q$      & $1,\ldots,K^*$          & component label of item $q$ \\
$g_q^{\dagger}$ & $1,\ldots,G^*$     & cluster label of item $q$ \\
$G^{*\,\dagger}$ & $\mathbb{N}$       & cluster truncation upper bound \\
$K^*$      & $\mathbb{N}$            & per-cluster component truncation upper bound \\
$K_g^{\dagger}$ & $\le K^*$          & nominal number of components in cluster $g$ \\
$\mu_l$    & $\R^d$                  & component mean ($l$ indexes a $(g,k)$ pair) \\
$\Sigma_l$ & $\R^{d\times d}$ PD     & component covariance \\
$\mu_g^{*\,\dagger}$ & $\R^d$        & cluster mean (centre of cluster $g$) \\
$\pi^{\dagger}$ & $\Delta^{G^*-1}$   & cluster mixture weights \\
$w_{\cdot|g}^{\dagger}$ & $\Delta^{K_g-1}$ & within-cluster component weights \\
$\alpha_0^{\dagger}$ & $\R_{>0}$     & cluster-level Dirichlet concentration \\
$\alpha_w^{\dagger}$ & $\R_{>0}$     & within-cluster Dirichlet concentration \\
$m_0$, $M_0^{\dagger}$ & $\R^d$, PD  & prior mean / covariance of cluster mean \\
$\kappa_0$ & $\R_{>0}$               & NIW component-mean precision factor \\
$\nu_0$, $S_0$ & $\R_{>0}$, PD       & NIW degrees of freedom / scale matrix \\
$\nu_W,\Lambda_W$ & $\R_{>0}$, PD    & Wishart hyperprior on $S_0$ \\
\midrule
$Y_{ij}^{(\ell)}$ & layer-specific  & observed response \\
$\eta_{ij}^{(\ell)}$ & $\R$         & LSIRM linear predictor \\
\bottomrule
\end{tabular}
\caption{Indices, dimensions, parameters, and observed quantities. New
v14 symbols are marked with $\dagger$.}
\label{tab:notation}
\end{table}

The global-to-local index map is denoted $q\mapsto(\ell(q),j(q))$,
inverse $j\mapsto q(\ell,j)$, exactly as in v13.

\section{LSIRM likelihood}
\label{sec:lsirm}

Layer-$\ell$ linear predictor:
\begin{equation}
\eta_{ij}^{(\ell)} \;=\; \alpha_i^{(\ell)} - \beta_j^{(\ell)}
   - \gamma_\ell \,\norm{a_i - b_j^{(\ell)}},
\qquad i=1,\ldots,N,\;\; j=1,\ldots,J_\ell.
\label{eq:eta}
\end{equation}

\subsection{Binary layer ($\ell=1$)}
\begin{equation}
Y_{ij}^{(1)}\mid\eta_{ij}^{(1)} \;\sim\; \mathrm{Bernoulli}\!\bigl(\sigma(\eta_{ij}^{(1)})\bigr),
\qquad
\sigma(x)=\bigl(1+e^{-x}\bigr)^{-1}.
\label{eq:bin}
\end{equation}

\subsection{Continuous layer ($\ell=2$, Student-$t$ via scale mixture)}
\begin{align}
\lambda_{ij}^{(2)} &\sim \mathrm{Gamma}(\nu_t/2,\nu_t/2),\\
Y_{ij}^{(2)}\mid\eta_{ij}^{(2)},\lambda_{ij}^{(2)},\sigma_0^2
   &\sim \mathcal{N}\!\Bigl(\eta_{ij}^{(2)},\;\sigma_0^2/\lambda_{ij}^{(2)}\Bigr),
\label{eq:cont}
\end{align}
which marginalises to $t_{\nu_t}(\eta_{ij}^{(2)},\sigma_0^2)$.

\subsection{Count layer ($\ell=3$, Negative-Binomial with item overdispersion)}
\begin{equation}
Y_{ij}^{(3)}\mid\eta_{ij}^{(3)},\kappa_j
   \;\sim\; \mathrm{NB}\!\bigl(\mu=\exp(\eta_{ij}^{(3)}),\;\phi=\kappa_j\bigr),
\qquad
\mathrm{Var}(Y) = \mu + \mu^2/\kappa_j.
\label{eq:nb}
\end{equation}

\subsection{Ordinal layer ($\ell=4$, LSGRM cumulative logit)}
With $H_\ell$ categories and increasing thresholds
$-\infty=\tau_0^{(\ell)}<\tau_1^{(\ell)}<\cdots<\tau_{H_\ell-1}^{(\ell)}<\tau_{H_\ell}^{(\ell)}=+\infty$,
\begin{equation}
\Pr(Y_{ij}^{(4)}\le h\mid\eta_{ij}^{(4)})
   = \sigma\!\bigl(\tau_h^{(4)} - \eta_{ij}^{(4)}\bigr),
\quad h=1,\ldots,H_\ell-1,
\label{eq:lsgrm}
\end{equation}
hence $\Pr(Y_{ij}^{(4)}=h)
   =\sigma(\tau_h^{(4)}-\eta_{ij}^{(4)})-\sigma(\tau_{h-1}^{(4)}-\eta_{ij}^{(4)})$.

\medskip
The likelihood factor is
\begin{equation}
p(Y\mid\Theta_{\text{LSIRM}})
 =\prod_{\ell=1}^{L}\prod_{i=1}^{N}\prod_{j=1}^{J_\ell}
   p_\ell\!\bigl(Y_{ij}^{(\ell)}\mid\eta_{ij}^{(\ell)},\cdot\bigr),
\label{eq:fulllik}
\end{equation}
identical to v13.

\section{Latent-position Mixture-of-Mixtures prior for items}
\label{sec:mom-prior}

\subsection{Global item position notation}

Stack the per-layer item positions into one global pool:
\begin{equation}
z_q \;\equiv\; b_{j(q)}^{(\ell(q))} \in \R^d,
\qquad q=1,\ldots,P,
\label{eq:zq}
\end{equation}
with $P=\sum_\ell J_\ell$. The MoM prior is placed on $\{z_q\}_{q=1}^P$.

\subsection{Two-level generative allocation}
\label{sec:gen}

For each item $q$, draw \emph{cluster} then \emph{component} then position:
\begin{align}
g_q\mid \pi               &\sim \Cat(\pi),
   \qquad \pi=(\pi_1,\ldots,\pi_{G^*})\in\Delta^{G^*-1}, \label{eq:gq}\\
c_q\mid g_q=g,\;w_{\cdot|g} &\sim \Cat(w_{\cdot|g}),
   \qquad w_{\cdot|g}=(w_{1|g},\ldots,w_{K_g|g})\in\Delta^{K_g-1}, \label{eq:cq}\\
z_q\mid c_q=l,\;\mu_l,\Sigma_l &\sim \mathcal{N}_d(\mu_l,\Sigma_l).\label{eq:zgivenc}
\end{align}
Component labels are scoped to a cluster: an index pair $(g,k)$ uniquely
identifies a component, and we will write $l\equiv(g,k)$ where convenient.
A cluster $g$ owns a private set of components
$\mathcal{L}_g=\{(g,1),\ldots,(g,K_g)\}$ with
$\mathcal{L}_g\cap\mathcal{L}_{g'}=\emptyset$ for $g\ne g'$.

The marginal distribution of $z_q$ is the two-level Gaussian mixture
\begin{equation}
p(z_q\mid \pi,\{w_{\cdot|g}\},\{\mu_l,\Sigma_l\})
 = \sum_{g=1}^{G^*}\pi_g
   \sum_{k=1}^{K_g} w_{k|g}\,
   \mathcal{N}_d(z_q;\mu_{(g,k)},\Sigma_{(g,k)}).
\label{eq:marg-z}
\end{equation}

\subsection{Component-parameter prior with cluster-mean centring}
\label{sec:nwprior}

For each component $l=(g,k)$, conditional on its cluster mean $\mu_g^*$,
\begin{align}
\Sigma_l \mid S_0,\nu_0 &\sim \IW(\nu_0,S_0), \label{eq:Sigprior}\\
\mu_l \mid g(l)=g,\,\mu_g^*,\,\Sigma_l,\kappa_0
   &\sim \mathcal{N}_d\!\Bigl(\mu_g^*,\;\Sigma_l/\kappa_0\Bigr).
\label{eq:muprior}
\end{align}
This is the standard Normal--Inverse-Wishart (NIW) prior, but with the
prior \emph{mean} replaced by the cluster centre $\mu_g^*$ rather than the
global anchor $m_0$. Conditional on $\mu_g^*$, $(\mu_l,\Sigma_l)$ is
exactly NIW$(\mu_g^*,\kappa_0,\nu_0,S_0)$, so all the NIW machinery from
v13 (Student-$t$ posterior predictive, marginal likelihood, etc.) applies
verbatim within each cluster, with $\mu_g^*$ playing the role of $m_0$.

\subsection{Cluster-mean prior with scale separation}
\label{sec:cmprior}

Cluster centres are independent across clusters with a wide Gaussian prior:
\begin{equation}
\mu_g^* \mid m_0,M_0 \;\sim\; \mathcal{N}_d(m_0,M_0),
\qquad g=1,\ldots,G^*.
\label{eq:must}
\end{equation}
We use a default isotropic $M_0 = m_{00}^2\,I_d$ with $m_0=0$, which makes
the prior rotation/reflection invariant about the origin
(see Sec.~\ref{sec:rotation}).

\paragraph{Scale separation.} The within-cluster mean spread is governed
by $\Sigma_l/\kappa_0$ in (\ref{eq:muprior}); marginalising
$\Sigma_l\sim\IW(\nu_0,S_0)$ gives
\begin{equation}
\mathbb{E}[\Sigma_l/\kappa_0]
 = \frac{1}{\kappa_0(\nu_0-d-1)} S_0,
\qquad \nu_0>d+1,
\label{eq:withinscale}
\end{equation}
whereas the between-cluster mean spread is $M_0=m_{00}^2 I_d$. The
\emph{cluster identifiability} condition is
\begin{equation}
\boxed{\;\;
m_{00}^2 \;\gg\; \frac{\tr(S_0)}{d\,\kappa_0(\nu_0-d-1)}.
\;\;}
\label{eq:scalesep}
\end{equation}
When this holds, the prior strongly distinguishes between
``components close together within a cluster'' (small displacement
$\mu_l-\mu_g^*$) and ``cluster centres spread widely'' (large
displacement $\mu_g^*-m_0$). In practice we recommend
$m_{00}^2/[\kappa_0(\nu_0-d-1)]^{-1}\tr(S_0)/d \;\ge\; 10^2$.
This is the prior mechanism that allows ``many close-together components
to form one non-Gaussian cluster'': within a cluster, the components
are pulled tightly toward $\mu_g^*$ at the within-scale; clusters
themselves are pulled loosely toward $m_0$ at the (much larger)
between-scale.

\subsection{Sparse Dirichlet weights at both levels}
\label{sec:dirichlet}

We use the Rousseau--Mengersen \emph{over-fitted finite mixture} setting
at both levels. Let $G^*$ be the cluster truncation and $K^*$ the per-cluster
component truncation (assume $K_g\equiv K^*$ for simplicity below; an
unequal-$K_g$ extension is a trivial relabelling).

\begin{align}
\pi          &\sim \Dir\!\bigl(\alpha_0/G^*,\;\ldots,\;\alpha_0/G^*\bigr), \label{eq:pi-prior}\\
w_{\cdot|g}  &\stackrel{\mathrm{iid}}{\sim} \Dir\!\bigl(\alpha_w/K^*,\;\ldots,\;\alpha_w/K^*\bigr),
\quad g=1,\ldots,G^*. \label{eq:w-prior}
\end{align}
Choosing $\alpha_0/G^* < d/2$ and $\alpha_w/K^* < d/2$ implements the
Rousseau--Mengersen sparse regime, in which posterior mass concentrates
on partitions of $\{1,\ldots,P\}$ with empty redundant components/clusters.
The number of \emph{occupied} clusters $G_+ = \#\{g:\sum_q \ind\{g_q=g\}>0\}$
and components $K_{+,g}=\#\{k:\sum_q \ind\{c_q=(g,k)\}>0\}$ are
posterior-identified summaries; the truncations $G^*,K^*$ are operational
upper bounds.

\subsection{Wishart hyperprior on $S_0$}
\label{sec:wishart-prior}

As in v13,
\begin{equation}
S_0 \sim \mathcal{W}_d(\nu_W,\Lambda_W),
\qquad \mathbb{E}[S_0]=\nu_W\Lambda_W,
\label{eq:wishprior}
\end{equation}
with $\nu_W>d-1$. This propagates uncertainty about the typical
component scale and avoids hand-tuning $S_0$.

\subsection{Cluster-identifiability mechanism: discussion}
\label{sec:cidisc}

In a flat Gaussian mixture, two components $l_1,l_2$ with similar
$(\mu_l,\Sigma_l)$ are statistically indistinguishable (label-switching).
In v14, the analogous label switching at the \emph{component} level is
benign and expected: components inside a cluster are exchangeable
density-basis atoms. What we want to identify is the \emph{cluster}
labelling, and this is achieved precisely when the prior makes within-cluster
component proximity overwhelmingly more probable than between-cluster
component proximity, i.e.\ (\ref{eq:scalesep}) holds.

If $m_{00}^2$ is too small (no scale separation), the model collapses
back to a one-level mixture: cluster labels become uninformative because
distinct clusters are no closer than two components in the same cluster.
If $m_{00}^2$ is too large, the prior may resist allowing a single cluster
to span an extended region of latent space, fragmenting an extended
cluster across several $g$'s. The default $m_{00}^2$ should be chosen so
that (\ref{eq:scalesep}) holds with a comfortable margin given the
expected support of $\{z_q\}$ in $\R^d$ (typically a few units in
LSIRM-Procrustes-anchored units).

\section{Priors for the LSIRM (non-clustering) parameters}
\label{sec:other-priors}

These are unchanged from v13:
\begin{align}
a_i              &\stackrel{\mathrm{iid}}{\sim} \mathcal{N}_d(0,\sigma_a^2 I_d), \label{eq:a-prior}\\
\alpha_i^{(\ell)}\mid \sigma_{\alpha,\ell}^2
                 &\stackrel{\mathrm{iid}}{\sim} \mathcal{N}(0,\sigma_{\alpha,\ell}^2),\label{eq:alpha-prior}\\
\sigma_{\alpha,\ell}^2
                 &\sim \IG(a_\alpha,b_\alpha),\label{eq:sigalpha-prior}\\
\beta_j^{(\ell)} &\stackrel{\mathrm{iid}}{\sim} \mathcal{N}(0,\sigma_\beta^2),\label{eq:beta-prior}\\
\gamma_\ell      &\stackrel{\mathrm{iid}}{\sim} \mathrm{LogNormal}(0,\sigma_\gamma^2),\label{eq:gamma-prior}\\
\sigma_0^2       &\sim \IG(a_\sigma,b_\sigma),\label{eq:sig0-prior}\\
\kappa_j         &\stackrel{\mathrm{iid}}{\sim} \mathrm{Gamma}(a_\kappa,b_\kappa),\label{eq:kappa-prior}\\
\tau_h^{(\ell)}  &= \tau_1^{(\ell)} + \sum_{r=2}^{h}\exp(\xi_r^{(\ell)}),
   \quad \tau_1^{(\ell)},\xi_r^{(\ell)}\sim\mathcal{N}(0,\sigma_\tau^2).\label{eq:tau-prior}
\end{align}
The item position $b_j^{(\ell)}$ does \emph{not} have an independent
prior; its prior is the MoM mixture (\ref{eq:marg-z}) via $z_{q(\ell,j)}=b_j^{(\ell)}$.

\section{Rotation and reflection invariance}
\label{sec:rotation}

The LSIRM linear predictor (\ref{eq:eta}) depends on $a_i,b_j^{(\ell)}$
only through $\norm{a_i - b_j^{(\ell)}}$ and so is invariant under any
rigid motion $T(x)=Rx+t$ with $R\in O(d)$, $t\in\R^d$, applied jointly
to all $a_i$, $b_j^{(\ell)}$. The prior on $a_i$ (centred Gaussian,
isotropic) and the priors on
$\mu_g^*$, $\mu_l$, $\Sigma_l$ must transform compatibly. Under
$T$:
\begin{align*}
\mu_g^*&\mapsto R\mu_g^*+t,\\
\mu_l  &\mapsto R\mu_l + t,\\
\Sigma_l&\mapsto R\Sigma_l R\trans.
\end{align*}
The component-level NIW (\ref{eq:Sigprior})--(\ref{eq:muprior}) is
$O(d)$-equivariant when $S_0\propto I_d$ (or $S_0$ itself is treated as a
Wishart variate as in (\ref{eq:wishprior}) with $\Lambda_W\propto I_d$).
The cluster-mean prior (\ref{eq:must}) is $O(d)$-equivariant when
$m_0=0$ and $M_0\propto I_d$. With these isotropic choices the joint
prior on $(\{a_i\},\{\mu_g^*\},\{\mu_l,\Sigma_l\})$, and hence the
posterior, is invariant under joint rotations/reflections about the
origin and translations in any direction, exactly matching v13's
discussion. Identifiability is restored at post-processing by
Procrustes alignment to a reference draw, which moves $\mu_g^*$ along
with $a_i,b_j^{(\ell)}$.

\section{Full joint posterior}
\label{sec:joint}

Let $\Theta_{\text{LSIRM}}=\{a_i,\alpha_i^{(\ell)},\beta_j^{(\ell)},
\gamma_\ell,\sigma_{\alpha,\ell}^2,\sigma_0^2,\lambda_{ij}^{(2)},
\kappa_j,\{\tau_h^{(\ell)}\}\}$ collect the LSIRM parameters and
$\Theta_{\text{MoM}}=\{g,c,\pi,\{w_{\cdot|g}\},\{\mu_g^*\},
\{\mu_l,\Sigma_l\},S_0\}$ collect the MoM parameters, where
$g=(g_1,\ldots,g_P)$ and $c=(c_1,\ldots,c_P)$.
Recall $z_q=b_{j(q)}^{(\ell(q))}$ is a function of $\Theta_{\text{LSIRM}}$.
The joint posterior factorises as
\begin{equation}
\begin{aligned}
&p(\Theta_{\text{LSIRM}},\Theta_{\text{MoM}}\mid Y) \;\propto\;
   p(Y\mid\Theta_{\text{LSIRM}})\\
&\quad\times \prod_{q=1}^{P}\Bigl[\pi_{g_q}\,w_{c_q\mid g_q}\,
   \mathcal{N}_d\!\bigl(z_q;\mu_{c_q},\Sigma_{c_q}\bigr)\Bigr]\\
&\quad\times \prod_{l}\Bigl[\mathcal{N}_d\!\bigl(\mu_l;\mu_{g(l)}^*,\Sigma_l/\kappa_0\bigr)\,
                          \IW(\Sigma_l;\nu_0,S_0)\Bigr]\\
&\quad\times \prod_{g=1}^{G^*}\mathcal{N}_d(\mu_g^*;m_0,M_0)\\
&\quad\times \Dir\!\bigl(\pi;\tfrac{\alpha_0}{G^*},\ldots\bigr)\,
             \prod_{g=1}^{G^*}\Dir\!\bigl(w_{\cdot|g};\tfrac{\alpha_w}{K^*},\ldots\bigr)\\
&\quad\times \mathcal{W}_d(S_0;\nu_W,\Lambda_W)\;\times\;p(\Theta_{\text{LSIRM}}^{(\text{non-}b)}).
\end{aligned}
\label{eq:joint}
\end{equation}
Equation (\ref{eq:joint}) is the operational target of the sampler. Note
the explicit \emph{joint} appearance of $g_q$ and $c_q$: the v14
likelihood--prior product treats them both as latent random variables
indexed at each $q$.

\section{NIW posterior quantities (cluster-conditional form)}
\label{sec:niw}

Within a single cluster $g$, conditional on $\mu_g^*$, the prior on each
component is the standard NIW, so the v13 NIW machinery applies with
$m_0\rightarrow\mu_g^*$.

\subsection{Component sufficient statistics}

Let $A_l = \{q:c_q=l\}$ be the items currently assigned to component $l$,
$n_l=|A_l|$. Define the component sufficient statistics
\begin{equation}
\bar z_l = \frac{1}{n_l}\sum_{q\in A_l} z_q,
\qquad
\mathcal{S}_l = \sum_{q\in A_l}(z_q-\bar z_l)(z_q-\bar z_l)\trans.
\label{eq:suffstats}
\end{equation}
With cluster-mean prior $\mu_g^*$ (where $g=g(l)$), the conjugate NIW
posterior parameters are
\begin{align}
\kappa_l^{(n)} &= \kappa_0+n_l, \label{eq:kn}\\
\nu_l^{(n)}    &= \nu_0+n_l, \label{eq:nn}\\
m_l^{(n)}      &= \frac{\kappa_0\,\mu_g^* + n_l\,\bar z_l}{\kappa_0+n_l}, \label{eq:mn}\\
S_l^{(n)}      &= S_0 + \mathcal{S}_l + \frac{\kappa_0 n_l}{\kappa_0+n_l}
                   (\bar z_l-\mu_g^*)(\bar z_l-\mu_g^*)\trans. \label{eq:Sn}
\end{align}

\subsection{Within-cluster collapsed Student-$t$ predictive}
\label{sec:tpred}

For an item $q$ with proposed component $l$, the v13 collapsed predictive
becomes (with $m_l^{(n)},S_l^{(n)},\kappa_l^{(n)},\nu_l^{(n)}$ all
conditioned on the current $\mu_g^*$):
\begin{equation}
p(z_q\mid c_q=l,\,z_{A_l},\,\mu_{g(l)}^*,S_0)
 = t_{\nu_l^{(n)}-d+1}\!\Bigl(z_q;\,m_l^{(n)},\;
       \tfrac{\kappa_l^{(n)}+1}{\kappa_l^{(n)}(\nu_l^{(n)}-d+1)}\,S_l^{(n)}\Bigr).
\label{eq:tpred}
\end{equation}
For an empty proposed component (fresh $l$ in cluster $g$), the
predictive uses the prior NIW with $\mu_g^*$ as prior mean:
\begin{equation}
p(z_q\mid c_q\!=\!l_{\mathrm{new}},\,g(l_{\mathrm{new}})=g,\,\mu_g^*,S_0)
 = t_{\nu_0-d+1}\!\Bigl(z_q;\,\mu_g^*,\;
   \tfrac{\kappa_0+1}{\kappa_0(\nu_0-d+1)}\,S_0\Bigr).
\label{eq:tpred-empty}
\end{equation}

\subsection{Within-cluster marginal likelihood of a component}

The marginal likelihood of the items in component $l$, with prior mean
$\mu_g^*$, is
\begin{equation}
\begin{aligned}
m_l(\mu_g^*) &\equiv p(z_{A_l}\mid \mu_g^*,S_0,\nu_0,\kappa_0)\\
 &= \pi^{-n_l d/2}
    \frac{\Gamma_d(\nu_l^{(n)}/2)}{\Gamma_d(\nu_0/2)}
    \frac{|S_0|^{\nu_0/2}}{|S_l^{(n)}|^{\nu_l^{(n)}/2}}
    \Bigl(\frac{\kappa_0}{\kappa_l^{(n)}}\Bigr)^{d/2},
\end{aligned}
\label{eq:marg-l}
\end{equation}
where $\Gamma_d$ is the multivariate gamma function. The
\emph{cluster-marginal likelihood} (after summing over component
allocations within the cluster) is
\begin{equation}
m_g(\{z_q:g_q=g\},\mu_g^*) =
   \sum_{\text{partitions }\mathcal{P}\text{ of }\{q:g_q=g\}\text{ into }\le K^*\text{ blocks}}
   p(\mathcal{P}\mid w_{\cdot|g},\alpha_w)\prod_{l\in\mathcal{P}} m_l(\mu_g^*),
\label{eq:marg-g}
\end{equation}
which is the marginal likelihood of a finite mixture and has no closed
form in $g$ alone. This is the obstruction that motivates the
conditional-on-$(\mu,\Sigma,\mu^*,w,\pi)$ split--merge variant of
Sec.~\ref{sec:splitmerge}.

\section{Collapsed component-label update $c_q\mid g_q$}
\label{sec:cupdate}

Conditional on $g_q=g$ and on the rest, $c_q$ is updated by ordinary
collapsed Gibbs \emph{within cluster $g$}, exactly as in v13 with
$\mu_g^*$ replacing $m_0$:
\begin{equation}
\Pr(c_q=l \mid g_q=g,\,\mathrm{rest})
\;\propto\;
\bigl(n_{l,-q} + \alpha_w/K^*\bigr)\,
p(z_q\mid c_q=l,z_{A_l\setminus\{q\}},\mu_g^*,S_0),
\label{eq:cupdate}
\end{equation}
for each $l\in\mathcal{L}_g$, where $n_{l,-q}=|A_l\setminus\{q\}|$ and
the predictive density is (\ref{eq:tpred}) (or (\ref{eq:tpred-empty})
for an empty $l$). Normalise over $l\in\mathcal{L}_g$ and sample.

This update preserves the within-cluster NIW collapsed structure exactly
and is the workhorse component move.

\section{Component-parameter posterior draws $(\mu_l,\Sigma_l)$}
\label{sec:musig-update}

For each occupied component $l$ (cluster $g=g(l)$), draw
\begin{align}
\Sigma_l \mid \cdot &\sim \IW\!\bigl(\nu_l^{(n)},\;S_l^{(n)}\bigr),\label{eq:Sig-draw}\\
\mu_l \mid \Sigma_l,\cdot &\sim \mathcal{N}_d\!\Bigl(m_l^{(n)},\;\Sigma_l/\kappa_l^{(n)}\Bigr).
\label{eq:mu-draw}
\end{align}
For each \emph{empty} (un-occupied) component, sample from the prior:
\begin{equation}
\Sigma_l\sim\IW(\nu_0,S_0),\qquad
\mu_l\mid\Sigma_l\sim\mathcal{N}_d(\mu_{g(l)}^*,\Sigma_l/\kappa_0).
\label{eq:emptydraw}
\end{equation}

\section{Cluster-mean posterior draw $\mu_g^*$}
\label{sec:mustar-update}

Given $\{\mu_l,\Sigma_l:g(l)=g\}$, the conditional for $\mu_g^*$ is
Gaussian--Gaussian conjugate. The likelihood contribution of the
components in cluster $g$ is
\begin{equation}
\prod_{l:g(l)=g}\mathcal{N}_d\!\bigl(\mu_l;\mu_g^*,\Sigma_l/\kappa_0\bigr)
\propto
\exp\!\Bigl\{-\tfrac{1}{2}\sum_{l:g(l)=g}
 \kappa_0(\mu_l-\mu_g^*)\trans\Sigma_l^{-1}(\mu_l-\mu_g^*)\Bigr\}.
\label{eq:mustar-lik}
\end{equation}
Combined with $\mathcal{N}_d(m_0,M_0)$ this yields
\begin{equation}
\mu_g^*\mid\cdot \sim \mathcal{N}_d\!\bigl(\hat m_g,\;\hat V_g\bigr),
\label{eq:mustar-post}
\end{equation}
with
\begin{align}
\hat V_g &=\Bigl(M_0^{-1}+\kappa_0\!\!\sum_{l:g(l)=g}\!\!\Sigma_l^{-1}\Bigr)^{-1},
   \label{eq:Vhat}\\
\hat m_g &=\hat V_g\Bigl(M_0^{-1}m_0
   +\kappa_0\!\!\sum_{l:g(l)=g}\!\!\Sigma_l^{-1}\mu_l\Bigr).
   \label{eq:mhat}
\end{align}
For \emph{empty} clusters $g$ ($\sum_q\ind\{g_q=g\}=0$ and consequently
no occupied components), draw $\mu_g^*\sim\mathcal{N}_d(m_0,M_0)$ from
the prior.

\section{Cluster-allocation update $g_q$}
\label{sec:gupdate}

Two valid Gibbs-style options:

\paragraph{Option A: collapsed $g_q$ (integrating out $c_q$).}
Marginalise $c_q$ analytically:
\begin{equation}
\Pr(g_q=g\mid\mathrm{rest}) \;\propto\;
\bigl(n_{g,-q} + \alpha_0/G^*\bigr)\,
\sum_{l\in\mathcal{L}_g}\bigl(n_{l,-q}+\alpha_w/K^*\bigr)\,
p(z_q\mid c_q=l,z_{A_l\setminus\{q\}},\mu_g^*,S_0),
\label{eq:goptA}
\end{equation}
where $n_{g,-q}=\sum_{r\ne q}\ind\{g_r=g\}$ and the predictive is
(\ref{eq:tpred})/(\ref{eq:tpred-empty}). After sampling $g_q$, draw $c_q$
from the within-cluster categorical (\ref{eq:cupdate}).

\paragraph{Option B (recommended): joint resampling of $(g_q,c_q)$.}
Treat the pair as a single categorical over
$\bigsqcup_g \mathcal{L}_g$:
\begin{equation}
\Pr(g_q=g,\,c_q=l\mid\mathrm{rest})\;\propto\;
(n_{g,-q}+\alpha_0/G^*)\,
(n_{l,-q}+\alpha_w/K^*)\,
p(z_q\mid c_q=l,z_{A_l\setminus\{q\}},\mu_{g(l)}^*,S_0).
\label{eq:goptB}
\end{equation}
Sample directly from this categorical of size $\sum_g K_g$. This is
exactly Option A followed by the conditional component draw, rolled into
one step; it preserves the NIW collapsed structure and is the recommended
default. We use Option B in Algorithm~\ref{alg:sweep}.

\section{Sparse Dirichlet weight updates}
\label{sec:weight-update}

Standard Dirichlet--Multinomial conjugacy.

\paragraph{Cluster weights.} With $n_g=\sum_q\ind\{g_q=g\}$,
\begin{equation}
\pi\mid g \;\sim\; \Dir\!\bigl(\alpha_0/G^*+n_1,\;\ldots,\;\alpha_0/G^*+n_{G^*}\bigr).
\label{eq:pi-update}
\end{equation}

\paragraph{Within-cluster component weights.} With $n_{(g,k)}=\sum_q\ind\{c_q=(g,k)\}$,
for each $g=1,\ldots,G^*$,
\begin{equation}
w_{\cdot|g}\mid c,g \;\sim\; \Dir\!\bigl(\alpha_w/K^*+n_{(g,1)},\;\ldots,\;\alpha_w/K^*+n_{(g,K^*)}\bigr).
\label{eq:w-update}
\end{equation}

\section{Wishart $S_0$ update}
\label{sec:wishart}

Carrying over the v13 conjugate update, with $L_+ = $ number of currently
existing components (occupied or not, but since empties contribute only
prior, we restrict to occupied $l$), summing across \emph{all clusters}:
\begin{equation}
S_0\mid\{\Sigma_l\} \;\sim\; \mathcal{W}_d\!\Bigl(\nu_W+L_+\nu_0,\;
   \bigl(\Lambda_W^{-1}+\sum_{l\text{ occ}}\Sigma_l^{-1}\bigr)^{-1}\Bigr).
\label{eq:S0-update}
\end{equation}

\section{Cluster-level split--merge target}
\label{sec:splitmerge}

\subsection{The cluster-marginal target}

A fully collapsed Jain--Neal split--merge on $g$ alone targets the
marginal posterior $p(g\mid Y,\Theta_{\text{LSIRM}}, S_0,\alpha_0,\alpha_w,
\nu_0,\kappa_0,m_0,M_0)$ obtained by integrating out
$(c,\{w_{\cdot|g}\},\{\mu_l,\Sigma_l\},\{\mu_g^*\},\pi)$:
\begin{equation}
p(g\mid\cdot) \;\propto\;
   \frac{\Gamma(\alpha_0)}{\Gamma(\alpha_0+P)}
   \prod_{g=1}^{G^*}\frac{\Gamma(n_g+\alpha_0/G^*)}{\Gamma(\alpha_0/G^*)}
   \cdot \prod_{g=1}^{G^*}\widetilde m_g(\{z_q:g_q=g\}),
\label{eq:marg-target}
\end{equation}
where $\widetilde m_g$ is the within-cluster marginal likelihood
$\int m_g(\{z\},\mu_g^*)\,\mathcal{N}_d(\mu_g^*;m_0,M_0)\,d\mu_g^*$
with $m_g$ as in (\ref{eq:marg-g}). This is intractable in two ways:
the sum over within-cluster partitions and the integral over $\mu_g^*$.

\subsection{Recommended: conditional-on-$(\mu,\Sigma,\mu^*,w,\pi)$ split--merge}
\label{sec:condSM}

Rather than collapse, we condition on the latest draws of
$(\{\mu_l,\Sigma_l\},\{\mu_g^*\},\{w_{\cdot|g}\},\pi)$ and target the joint
\begin{equation}
\pi(g,c \mid z,\{\mu_l,\Sigma_l\},\{\mu_g^*\},\{w_{\cdot|g}\},\pi)
 \;\propto\;
\prod_{q=1}^{P}\bigl[\pi_{g_q}\,w_{c_q|g_q}\,
   \mathcal{N}_d(z_q;\mu_{c_q},\Sigma_{c_q})\bigr].
\label{eq:joint-cond}
\end{equation}
This density is closed-form in $(g,c)$, and any Metropolis--Hastings move
that proposes new $(g',c')$ and accepts with the (\ref{eq:joint-cond})
ratio is a valid Gibbs-within-MH step targeting the correct full posterior
(\ref{eq:joint}). This is the v14 recommended split--merge target.

We also note that, when an MH move is rejected, the chain remains at
$(g,c)$ with the marginal-stationary distribution maintained automatically;
intractability of (\ref{eq:marg-target}) is bypassed.

\section{Metropolis--Hastings updates for LSIRM parameters}
\label{sec:mh}

These are unchanged from v13. The only place the MoM prior enters is in
the $b_j^{(\ell)}$ update, where the mixture-prior term uses the current
$(\mu_{c_q},\Sigma_{c_q})$ for the appropriate component label $c_q$ of
$q=q(\ell,j)$.

\subsection{Respondent positions $a_i$}
\begin{equation}
a_i^{\mathrm{prop}} \sim \mathcal{N}_d(a_i,\Sigma_a^{\mathrm{prop}}),
\quad
\log r_a = \sum_{\ell,j}\log
   \frac{p_\ell(Y_{ij}^{(\ell)}\mid\eta_{ij}^{(\ell)\,\mathrm{prop}},\cdot)}
        {p_\ell(Y_{ij}^{(\ell)}\mid\eta_{ij}^{(\ell)},\cdot)}
   - \frac{1}{2\sigma_a^2}\bigl(\norm{a_i^{\mathrm{prop}}}^2-\norm{a_i}^2\bigr).
\label{eq:a-mh}
\end{equation}

\subsection{Respondent abilities $\alpha_i^{(\ell)}$}
\begin{equation}
\alpha_i^{(\ell)\mathrm{prop}}\sim\mathcal{N}(\alpha_i^{(\ell)},(\sigma_\alpha^{\mathrm{prop}})^2),
\quad
\log r_\alpha = \sum_j\log\frac{p_\ell(\cdot\mid\eta^{\mathrm{prop}})}{p_\ell(\cdot\mid\eta)}
   - \frac{(\alpha_i^{(\ell)\mathrm{prop}})^2 - (\alpha_i^{(\ell)})^2}{2\sigma_{\alpha,\ell}^2}.
\label{eq:alpha-mh}
\end{equation}

\subsection{Item easinesses $\beta_j^{(\ell)}$ and ordinal thresholds}
\begin{equation}
\beta_j^{(\ell)\mathrm{prop}}\sim\mathcal{N}(\beta_j^{(\ell)},(\sigma_\beta^{\mathrm{prop}})^2),
\quad
\log r_\beta = \sum_i\log\frac{p_\ell(\cdot\mid\eta^{\mathrm{prop}})}{p_\ell(\cdot\mid\eta)}
   - \frac{(\beta_j^{(\ell)\mathrm{prop}})^2 - (\beta_j^{(\ell)})^2}{2\sigma_\beta^2}.
\label{eq:beta-mh}
\end{equation}
Ordinal thresholds $\tau_h^{(4)}$ are updated through their unconstrained
parametrisation $(\tau_1^{(4)},\xi_2^{(4)},\ldots)$ via random-walk MH on
each component.

\subsection{Item positions $b_j^{(\ell)}$ with mixture prior}

This is the only LSIRM update where the MoM prior appears. Let
$q=q(\ell,j)$ and $l=c_q$ (the current component label). With optional
inflation $\alpha\ge 1$ (a sampler-level decoupling factor; see v13),
the proposal is
\begin{equation}
b_j^{(\ell)\mathrm{prop}} \sim \mathcal{N}_d\!\bigl(b_j^{(\ell)},\;\alpha\,\Sigma_b^{\mathrm{prop}}\bigr).
\label{eq:b-prop}
\end{equation}
The MH log-ratio is
\begin{equation}
\begin{aligned}
\log r_b
&= \sum_i\log\frac{p_\ell(Y_{ij}^{(\ell)}\mid\eta_{ij}^{(\ell)\,\mathrm{prop}},\cdot)}
                  {p_\ell(Y_{ij}^{(\ell)}\mid\eta_{ij}^{(\ell)},\cdot)}\\
&\quad
   -\tfrac12\bigl[
   (b_j^{(\ell)\mathrm{prop}}-\mu_{c_q})\trans\Sigma_{c_q}^{-1}(b_j^{(\ell)\mathrm{prop}}-\mu_{c_q})
  -(b_j^{(\ell)}-\mu_{c_q})\trans\Sigma_{c_q}^{-1}(b_j^{(\ell)}-\mu_{c_q})\bigr],
\end{aligned}
\label{eq:b-mh}
\end{equation}
identical in form to v13's eq.~(31) but with $\mu_{c_q},\Sigma_{c_q}$
drawn from the v14 two-level prior.

\paragraph{The inflation factor $\alpha$.} Setting $\alpha>1$ widens the
proposal beyond the prior scale and encourages the LSIRM geometry to be
explored independently of the cluster identity, which empirically reduces
posterior coupling between $(c_q,g_q)$ and $b_j^{(\ell)}$. The MH ratio
(\ref{eq:b-mh}) does not depend on $\alpha$ since the proposal is
symmetric.

\subsection{Other LSIRM MH updates}
$\gamma_\ell$, $\sigma_0^2$, $\lambda_{ij}^{(2)}$, $\kappa_j$ are
updated by their v13 MH (or, for $\lambda_{ij}^{(2)}$, by direct
Gamma conjugate Gibbs given the Student-$t$ scale-mixture
representation):
\begin{equation}
\lambda_{ij}^{(2)}\mid Y_{ij}^{(2)},\eta_{ij}^{(2)},\sigma_0^2 \sim
\mathrm{Gamma}\!\Bigl(\tfrac{\nu_t+1}{2},
\tfrac{\nu_t}{2}+\tfrac{(Y_{ij}^{(2)}-\eta_{ij}^{(2)})^2}{2\sigma_0^2}\Bigr).
\label{eq:lambda-update}
\end{equation}

\section{Conjugate variance updates}
\label{sec:varupd}

Carrying over the v13 inverse-gamma conjugate update for the
ability variance:
\begin{equation}
\sigma_{\alpha,\ell}^2\mid\{\alpha_i^{(\ell)}\} \;\sim\;
\IG\!\Bigl(a_\alpha+\tfrac{N}{2},\;
   b_\alpha+\tfrac12\sum_i (\alpha_i^{(\ell)})^2\Bigr).
\label{eq:sigalpha-update}
\end{equation}
Likewise, with the standard inverse-gamma--normal conjugate structure,
\begin{equation}
\sigma_0^2\mid\cdots \;\sim\;
\IG\!\Bigl(a_\sigma+\tfrac{NJ_2}{2},\;
   b_\sigma+\tfrac12\sum_{i,j}\lambda_{ij}^{(2)}\bigl(Y_{ij}^{(2)}-\eta_{ij}^{(2)}\bigr)^2\Bigr).
\label{eq:sig0-update}
\end{equation}

\section{Jain--Neal cluster-level split--merge}
\label{sec:jainneal}

We give a complete derivation of the cluster-level split--merge
sub-sampler, targeting (\ref{eq:joint-cond}) (the conditional-on
$(\{\mu_l,\Sigma_l\},\{\mu_g^*\},\{w_{\cdot|g}\},\pi)$ joint of $(g,c)$).
This is run $M_{SM}$ times per sweep.

\subsection{Anchor selection}

Sample anchors $u\ne v$ uniformly from $\{1,\ldots,P\}$.
Let $A=\{q:g_q\in\{g_u,g_v\}\}$ be the set of items in the union of the
two anchor clusters; let $A_*=A\setminus\{u,v\}$.
\begin{itemize}
\item If $g_u=g_v=A_0$: \emph{split} attempt. Choose a fresh empty
   cluster label $g_b\in\{1,\ldots,G^*\}\setminus\{\text{occupied}\}$
   (uniformly among empties; if none exists, the move is null and the
   sweep records a no-op). Initialise $g_u\leftarrow A_0\equiv g_a$,
   $g_v\leftarrow g_b$.
\item If $g_u\ne g_v$: \emph{merge} attempt. Let
   $g_a=g_u$, $g_b=g_v$. The merge proposal sends every $q\in A$ to
   $g_a$.
\end{itemize}

\subsection{Restricted Gibbs scan (split case)}

For each $q\in A_*$ (in random order), draw $g_q\in\{g_a,g_b\}$ from the
restricted conditional under (\ref{eq:joint-cond}). Define for each
candidate cluster $g\in\{g_a,g_b\}$ a within-cluster predictive (a finite
Gaussian mixture using the current $\{\mu_l,\Sigma_l\}_{l\in\mathcal{L}_g}$
and weights $\{w_{l|g}\}$):
\begin{equation}
f_g(z_q) \;\equiv\; \sum_{l\in\mathcal{L}_g} w_{l|g}\,
   \mathcal{N}_d(z_q;\mu_l,\Sigma_l).
\label{eq:fg}
\end{equation}
Then the restricted conditional is
\begin{equation}
\Pr_{\mathrm{prop}}(g_q=g\mid g_{A_*\setminus\{q\}})
 \;\propto\;
 (n_{g,-q}^{(A)} + \alpha_0/G^*)\,f_g(z_q),
 \quad g\in\{g_a,g_b\},
\label{eq:restgibbs}
\end{equation}
where $n_{g,-q}^{(A)}$ counts items in $A$ currently assigned to $g$
\emph{excluding} $q$. The restricted scan produces a launch state
$(g^{\mathrm{launch}},c^{\mathrm{launch}})$ after $T$ sweeps over $A_*$
(typically $T=5$). For $q\in A_*$, the corresponding component label is
also resampled: given the proposed $g_q$, draw
\begin{equation}
\Pr(c_q=l\mid g_q=g)\;\propto\;w_{l|g}\,\mathcal{N}_d(z_q;\mu_l,\Sigma_l),
\quad l\in\mathcal{L}_g.
\label{eq:csplitc}
\end{equation}
The forward proposal density $q_{\mathrm{fwd}}(g',c'\mid g,c)$ is the
product over the final restricted-scan probabilities (\ref{eq:restgibbs})
and (\ref{eq:csplitc}).

\paragraph{Component re-attachment under cluster split.}
When cluster $A_0$ splits into $g_a$ and $g_b$, the component pool
$\mathcal{L}_{A_0}$ must be repartitioned. We re-label: each $l\in
\mathcal{L}_{A_0}$ with at least one item in the new $g_a$ keeps a
position in $\mathcal{L}_{g_a}$; if it also has items in $g_b$, its label
is duplicated by a fresh slot in $\mathcal{L}_{g_b}$ initialised to the
same $(\mu_l,\Sigma_l)$ (with $w_{l|g_b}$ initialised proportional to
the count of its items in $g_b$). This step is deterministic given the
final $g$-labels of the items, since component labels are recomputed
from the within-cluster Gibbs (\ref{eq:csplitc}).

\subsection{Merge case}

The merge proposal is deterministic: set $g_q=g_a$ for all $q\in A$,
and reuse component labels $c_q$ unchanged (since the component pools
$\mathcal{L}_{g_a}$ and $\mathcal{L}_{g_b}$ are now combined into
$\mathcal{L}_{g_a}$, this is a relabelling). The reverse-split proposal
density is computed by running the same restricted Gibbs scan
\emph{from} the merged state \emph{back to} the original split state,
using the same anchors $u,v$ and the same $T$ sweeps; this is standard
Jain--Neal.

\subsection{Acceptance ratios}

Let $(g,c)$ be the current state and $(g',c')$ the proposal. The MH
acceptance probability is $\min\{1,r\}$, with
\begin{equation}
r = \frac{\pi(g',c'\mid z,\cdot)}{\pi(g,c\mid z,\cdot)}\cdot
    \frac{q_{\mathrm{rev}}(g,c\mid g',c')}{q_{\mathrm{fwd}}(g',c'\mid g,c)}.
\label{eq:r-mh}
\end{equation}

\paragraph{Split acceptance.} Substituting (\ref{eq:joint-cond}):
\begin{equation}
\begin{aligned}
\log r_{\mathrm{split}}
 &= \log\frac{\Gamma(\alpha_0/G^*+n_{g_a}^{\prime})\,\Gamma(\alpha_0/G^*+n_{g_b}^{\prime})}
              {\Gamma(\alpha_0/G^*+n_{A_0})\,\Gamma(\alpha_0/G^*)}\\
 &\quad + \sum_{q\in A}\Bigl[
    \log\pi_{g'_q}+\log w_{c'_q|g'_q}+\log\mathcal{N}_d(z_q;\mu_{c'_q},\Sigma_{c'_q})\\
 &\hspace{45pt}-\log\pi_{g_q}-\log w_{c_q|g_q}-\log\mathcal{N}_d(z_q;\mu_{c_q},\Sigma_{c_q})\Bigr]\\
 &\quad +\log q_{\mathrm{rev}}(\text{merge}\mid\text{split})
        -\log q_{\mathrm{fwd}}(\text{split}\mid\text{merge}).
\end{aligned}
\label{eq:rsplit}
\end{equation}
Here $n_{g_a}',n_{g_b}'$ are the post-proposal cluster sizes within $A$,
and $n_{A_0}=n_{g_a}'+n_{g_b}'$.
The merge-as-reverse proposal is deterministic, so
$\log q_{\mathrm{rev}}=0$ and the only stochastic term is the forward
restricted scan, $\log q_{\mathrm{fwd}}$.

\paragraph{Merge acceptance.} The reverse split is the stochastic move,
giving
\begin{equation}
\begin{aligned}
\log r_{\mathrm{merge}}
 &= \log\frac{\Gamma(\alpha_0/G^*+n_{g_a}^{\prime\prime})\,\Gamma(\alpha_0/G^*)}
              {\Gamma(\alpha_0/G^*+n_{g_a})\,\Gamma(\alpha_0/G^*+n_{g_b})}\\
 &\quad + \sum_{q\in A}\Bigl[\log\pi_{g'_q}+\log w_{c'_q|g'_q}+\log\mathcal{N}_d(z_q;\mu_{c'_q},\Sigma_{c'_q})\\
 &\hspace{45pt}-\log\pi_{g_q}-\log w_{c_q|g_q}-\log\mathcal{N}_d(z_q;\mu_{c_q},\Sigma_{c_q})\Bigr]\\
 &\quad +\log q_{\mathrm{rev}}(\text{split}\mid\text{merge})
        -\log q_{\mathrm{fwd}}(\text{merge}\mid\text{split}),
\end{aligned}
\label{eq:rmerge}
\end{equation}
with $n_{g_a}^{\prime\prime}=n_{g_a}+n_{g_b}$ and
$\log q_{\mathrm{fwd}}(\text{merge}\mid\text{split})=0$ because merge is
deterministic. Equations (\ref{eq:rsplit}) and (\ref{eq:rmerge}) mirror
v13's eqs.~(40) and~(43), now at the cluster level.

\subsection{Validity}

The move targets the conditional posterior
$\pi(g,c\mid z,\{\mu_l,\Sigma_l\},\{\mu_g^*\},\{w_{\cdot|g}\},\pi)$,
which is closed-form. Combined with the Gibbs scans for the conditioned
quantities (Secs.~\ref{sec:musig-update}, \ref{sec:mustar-update},
\ref{sec:weight-update}, \ref{sec:wishart}), this yields a valid
Metropolis-within-Gibbs sampler with stationary distribution
(\ref{eq:joint}). The fully-collapsed alternative (\ref{eq:marg-target})
would require a Monte Carlo estimator of $\widetilde m_g$ (e.g.\ a Chib
estimator or harmonic-mean variant), which we do not pursue here.

\section{Posterior cluster summaries}
\label{sec:summaries}

\subsection{Cluster-level posterior similarity}

The primary clustering summary is the cluster-level posterior similarity
matrix
\begin{equation}
C_{qr}^{\mathrm{cluster}} \;=\; \Pr(g_q=g_r\mid Y),
\qquad q,r=1,\ldots,P,
\label{eq:Ccluster}
\end{equation}
estimated by the Monte Carlo average
$\widehat C_{qr}^{\mathrm{cluster}}=\frac{1}{M}\sum_{m=1}^M\ind\{g_q^{(m)}=g_r^{(m)}\}$
over post-burn-in iterations. Point clusterings are obtained by
minimising Binder's loss or VI loss against $\widehat C^{\mathrm{cluster}}$,
exactly as in v13.

\subsection{Component-level posterior similarity}

The secondary, finer-grained summary is
\begin{equation}
C_{qr}^{\mathrm{component}} \;=\; \Pr(c_q=c_r\mid Y),
\label{eq:Ccomponent}
\end{equation}
estimated identically. Note that
$\ind\{c_q=c_r\}\Rightarrow\ind\{g_q=g_r\}$ (components are nested in
clusters), so $\widehat C^{\mathrm{component}}\le \widehat C^{\mathrm{cluster}}$
elementwise.

\subsection{Occupancy traces}

Track the number of occupied clusters
\begin{equation}
G_+^{(m)} = \#\{g:\sum_q\ind\{g_q^{(m)}=g\}>0\},
\label{eq:Gplus}
\end{equation}
and the per-cluster component occupancies
\begin{equation}
K_{+,g}^{(m)} = \#\{k:\sum_q\ind\{c_q^{(m)}=(g,k)\}>0\},
\label{eq:Kplus}
\end{equation}
across iterations $m$. Posterior modes / quantiles of $G_+,K_{+,g}$
provide a principled way to read off the effective number of clusters
and the within-cluster density complexity, with the truncations
$G^*,K^*$ acting only as upper bounds.

\section{Full MCMC sweep}
\label{sec:sweep}

We give the algorithmic specification in two algorithm boxes:
the LSIRM/MoM Gibbs sweep (Algorithm~\ref{alg:sweep}) and the cluster-level
split--merge sub-routine (Algorithm~\ref{alg:sm}).

\begin{algorithm}[H]
\caption{One full sweep of the v14 MCMC sampler.}
\label{alg:sweep}
\begin{algorithmic}[1]
\State \textbf{LSIRM block:}
\For{$i=1,\ldots,N$}
   \State Update $a_i$ via random-walk MH (eq.~\ref{eq:a-mh}).
\EndFor
\For{$\ell=1,\ldots,L$}
   \For{$i=1,\ldots,N$} \State Update $\alpha_i^{(\ell)}$ via MH (eq.~\ref{eq:alpha-mh}). \EndFor
   \For{$j=1,\ldots,J_\ell$} \State Update $\beta_j^{(\ell)}$ via MH (eq.~\ref{eq:beta-mh}). \EndFor
   \For{$j=1,\ldots,J_\ell$}
      \State $q\gets q(\ell,j)$;
             update $b_j^{(\ell)}$ via MH (eq.~\ref{eq:b-mh}) using $(\mu_{c_q},\Sigma_{c_q})$.
   \EndFor
   \State Update $\gamma_\ell$ via MH; if $\ell=4$ update ordinal thresholds via MH.
   \State Update $\sigma_{\alpha,\ell}^2$ via Gibbs (eq.~\ref{eq:sigalpha-update}).
\EndFor
\State Update $\sigma_0^2$ via Gibbs (eq.~\ref{eq:sig0-update}).
\For{$i,j$ continuous-layer pairs}
   \State Update $\lambda_{ij}^{(2)}$ via Gibbs (eq.~\ref{eq:lambda-update}).
\EndFor
\For{$j=1,\ldots,J_3$} \State Update $\kappa_j$ via MH. \EndFor
\State \textbf{Build global pool:} for each $q$, set $z_q\gets b_{j(q)}^{(\ell(q))}$.
\State \textbf{MoM allocation block:}
\For{$q=1,\ldots,P$}
   \State Sample $(g_q,c_q)$ jointly from (\ref{eq:goptB}). \Comment{Option B, recommended}
\EndFor
\State \textbf{Component-parameter block:}
\For{$l$ occupied}
   \State Compute $(\kappa_l^{(n)},\nu_l^{(n)},m_l^{(n)},S_l^{(n)})$ from (\ref{eq:kn})--(\ref{eq:Sn}) using $\mu_{g(l)}^*$.
   \State Draw $\Sigma_l$ from (\ref{eq:Sig-draw}); draw $\mu_l$ from (\ref{eq:mu-draw}).
\EndFor
\State For empty $l$, sample from prior (\ref{eq:emptydraw}).
\State \textbf{Cluster-mean block:}
\For{$g$ with at least one occupied component}
   \State Compute $(\hat m_g,\hat V_g)$ from (\ref{eq:Vhat})--(\ref{eq:mhat}); draw $\mu_g^*$ from (\ref{eq:mustar-post}).
\EndFor
\State For empty $g$, draw $\mu_g^*\sim\mathcal{N}_d(m_0,M_0)$.
\State \textbf{Sparse Dirichlet block:}
\State Draw $\pi$ from (\ref{eq:pi-update}); for each $g$, draw $w_{\cdot|g}$ from (\ref{eq:w-update}).
\State \textbf{Wishart block:} draw $S_0$ from (\ref{eq:S0-update}).
\State \textbf{Split--merge block:}
\For{$m=1,\ldots,M_{SM}$}
   \State Run Algorithm~\ref{alg:sm}.
\EndFor
\State \textbf{Storage:} store $g, c, K_+, G_+, z, \Theta_{\text{LSIRM}}$.
\State \textbf{Online update:} update running sum
   $\widehat C_{qr}^{\mathrm{cluster}}\!+\!\!=\ind\{g_q=g_r\}$;
   similarly for $\widehat C^{\mathrm{component}}$.
\end{algorithmic}
\end{algorithm}

\begin{algorithm}[H]
\caption{One Jain--Neal cluster-level split--merge step.}
\label{alg:sm}
\begin{algorithmic}[1]
\State Sample anchors $u\ne v$ uniformly from $\{1,\ldots,P\}$.
\State $A\gets\{q:g_q\in\{g_u,g_v\}\}$, $A_*\gets A\setminus\{u,v\}$.
\If{$g_u=g_v$} \Comment{split attempt}
   \State Pick fresh empty label $g_b$; $g_a\gets g_u$.
   \State Initialise $g_u^{\mathrm{launch}}=g_a$, $g_v^{\mathrm{launch}}=g_b$;
          for $q\in A_*$ pick $g_q^{\mathrm{launch}}\in\{g_a,g_b\}$ uniformly.
   \For{$t=1,\ldots,T$}
      \For{$q\in A_*$ in random order}
         \State Draw $g_q^{\mathrm{launch}}$ from the restricted conditional (\ref{eq:restgibbs}).
         \State Draw $c_q^{\mathrm{launch}}$ from (\ref{eq:csplitc}) given $g_q^{\mathrm{launch}}$.
      \EndFor
   \EndFor
   \State $(g',c')\gets(g^{\mathrm{launch}},c^{\mathrm{launch}})$;
          compute $\log q_{\mathrm{fwd}}$ as the product of (\ref{eq:restgibbs}) and (\ref{eq:csplitc})
          probabilities of the realised draws on the \emph{final} scan.
   \State Compute $\log r_{\mathrm{split}}$ via (\ref{eq:rsplit}).
   \State Accept $(g',c')$ with probability $\min\{1,e^{\log r_{\mathrm{split}}}\}$.
\Else \Comment{merge attempt}
   \State $g_a\gets g_u$, $g_b\gets g_v$.
   \State Set $g'_q=g_a$ for all $q\in A$; relabel components in
          $\mathcal{L}_{g_b}$ as fresh slots in $\mathcal{L}_{g_a}$;
          set $c'_q$ accordingly.
   \State Run a hypothetical reverse split: starting from a launch state
          with $g_u=g_a,g_v=g_b$ and the original split labels for
          $q\in A_*$, run $T$ restricted-Gibbs scans and compute
          $\log q_{\mathrm{rev}}$ as the probability of realising the
          original split state on the final scan.
   \State Compute $\log r_{\mathrm{merge}}$ via (\ref{eq:rmerge}).
   \State Accept $(g',c')$ with probability $\min\{1,e^{\log r_{\mathrm{merge}}}\}$.
\EndIf
\end{algorithmic}
\end{algorithm}

\section{Closing remarks}

The v14 MoM prior delivers two practical benefits over v13:
\begin{enumerate}[label=(\roman*)]
\item \emph{Non-Gaussian clusters by construction.} A cluster is a
   finite mixture of Gaussians, so elongated, skewed, or multi-modal
   clusters need no model adjustment.
\item \emph{Identifiability of the cluster level.} The scale separation
   (\ref{eq:scalesep}) gives the cluster labels $g_q$ an interpretation
   that flat mixtures lack: components close in $\R^d$ form one cluster,
   components far apart cannot.
\end{enumerate}
The cost is that the cluster-marginal likelihood (\ref{eq:marg-g}) is
intractable; the recommended remedy is the conditional-on-$(\mu,\Sigma,
\mu^*,w,\pi)$ Jain--Neal variant of Sec.~\ref{sec:condSM}, which
sidesteps this intractability while still targeting the correct full
posterior. With sparse Dirichlet weights at both levels, redundant
components within a cluster and redundant clusters die out automatically;
the truncations $G^*$ and $K^*$ play only the role of operational upper
bounds, and the posterior summaries $G_+$ and $K_{+,g}$ are the
substantive quantities to report.

All other aspects of the model --- the LSIRM linear predictor,
the four-layer likelihood, the respondent-side priors, the rotation/
reflection invariance, the optional inflation factor $\alpha\ge 1$ in
the $b_j$ MH update, and the inverse-gamma conjugate updates --- are
inherited verbatim from v13.

\end{document}